% Created 2026-11-26 jeu. 18:10
% Intended LaTeX compiler: pdflatex
\documentclass[paper=a4, fontsize=12pt, DIV=13, english, dvipsnames]{scrartcl}
\usepackage[utf8]{inputenc}
\usepackage[T1]{fontenc}
\usepackage{graphicx}
\usepackage{longtable}
\usepackage{wrapfig}
\usepackage{rotating}
\usepackage[normalem]{ulem}
\usepackage{amsmath}
\usepackage{amssymb}
\usepackage{hyperref}
\usepackage{amsfonts}
\usepackage{bookmark}
\usepackage{xcolor}
\definecolor{myblue}{RGB}{0,101,165}
\hypersetup{colorlinks=true, allcolors=myblue}
\definecolor{bg}{rgb}{0.95,0.95,0.95}
\definecolor{darkgreen}{RGB}{0,150,0}
\usepackage{booktabs}
\usepackage{colortbl}
\addtokomafont{title}{\color{myblue}}
\usepackage[modulo]{lineno}
\usepackage{float}
\pdfstringdefDisableCommands{\def\and{and }}
\setcapindent{0em}
\author{Ghislain Vieilledent$^{\star, 1, 2, 3}$ \and Laura Lagourgue$^{4}$ \and Marjolaine David$^{4}$ \and Adrien Bertaud$^{4}$ \and Oriane Bruyère$^{4}$ \and Philippe Birnbaum$^{1, 2, 3}$ \and Thomas Ibanez$^{1, 5}$ \and Nathan Eltabet$^{1, 5}$ \and David Bruy$^{1, 5}$ \and Patrick Heuret$^{1, 5, 6}$ \and Fabien Albouy$^{4}$}
\date{}
\title{Validation of Global Forest Cover Products and Analysis of Deforestation, Forest Regrowth and Forest Fragmentation in New Caledonia (2005–2025)}
\hypersetup{
 pdfauthor={Ghislain Vieilledent$^{\star, 1, 2, 3}$ \and Laura Lagourgue$^{4}$ \and Marjolaine David$^{4}$ \and Adrien Bertaud$^{4}$ \and Oriane Bruyère$^{4}$ \and Philippe Birnbaum$^{1, 2, 3}$ \and Thomas Ibanez$^{1, 5}$ \and Nathan Eltabet$^{1, 5}$ \and David Bruy$^{1, 5}$ \and Patrick Heuret$^{1, 5, 6}$ \and Fabien Albouy$^{4}$},
 pdftitle={Validation of Global Forest Cover Products and Analysis of Deforestation, Forest Regrowth and Forest Fragmentation in New Caledonia (2005–2025)},
 pdfkeywords={},
 pdfsubject={},
 pdfcreator={Emacs 30.2 (Org mode 9.7.11)}, 
 pdflang={English}}
\usepackage{calc}
\newlength{\cslhangindent}
\setlength{\cslhangindent}{1.5em}
\newlength{\csllabelsep}
\setlength{\csllabelsep}{0.6em}
\newlength{\csllabelwidth}
\setlength{\csllabelwidth}{0.45em * 0}
\newenvironment{cslbibliography}[2] % 1st arg. is hanging-indent, 2nd entry spacing.
 {% By default, paragraphs are not indented.
  \setlength{\parindent}{0pt}
  % Hanging indent is turned on when first argument is 1.
  \ifodd #1
  \let\oldpar\par
  \def\par{\hangindent=\cslhangindent\oldpar}
  \fi
  % Set entry spacing based on the second argument.
  \setlength{\parskip}{\parskip +  #2\baselineskip}
 }%
 {}
\newcommand{\cslblock}[1]{#1\hfill\break}
\newcommand{\cslleftmargin}[1]{\parbox[t]{\csllabelsep + \csllabelwidth}{#1}}
\newcommand{\cslrightinline}[1]
  {\parbox[t]{\linewidth - \csllabelsep - \csllabelwidth}{#1}\break}
\newcommand{\cslindent}[1]{\hspace{\cslhangindent}#1}
\newcommand{\cslbibitem}[2]
  {\leavevmode\vadjust pre{\hypertarget{citeproc_bib_item_#1}{}}#2}
\makeatletter
\newcommand{\cslcitation}[2]
 {\protect\hyper@linkstart{cite}{citeproc_bib_item_#1}#2\hyper@linkend}
\makeatother\begin{document}

\maketitle
\vspace{-1.5cm}
\begin{center}
{\small
$^{1}$AMAP, Univ Montpellier, CIRAD, CNRS, INRAE, IRD, Montpellier, FRANCE\\
$^{2}$CIRAD, UMR AMAP, F-34398 Montpellier, FRANCE\\
$^{3}$IAC, Nouméa, NEW CALEDONIA\\
$^{4}$OEIL, Nouméa, NEW CALEDONIA\\
$^{5}$IRD, UMR AMAP, Nouméa, NEW CALEDONIA\\
$^{6}$INRAE, UMR AMAP, Montpellier, FRANCE\\
$^{\star}$Correspondance: \href{mailto:ghislain.vieilledent@cirad.fr}{\texttt{ghislain.vieilledent@cirad.fr}}}
\end{center}
\linenumbers
\vspace{0.10cm}
\enlargethispage{\baselineskip}

\begin{abstract}
New Caledonia has been classified as a hotspot of biodiversity.
\par\vskip\baselineskip\noindent
\small \textbf{Keywords}: afforestation, biodiversity hotspot, deforestation, forest cover, land use change, photo-interpretation, regrowth, tropical forests, New-Caledonia, uncertainty.\\
\textbf{Running title}: Forest cover change in New Caledonia (2005--2025).
\end{abstract}

\newpage
\section{Introduction}
\label{sec:org61f91ef}

\newpage
\section{Materials and Methods}
\label{sec:orgd815640}

\subsection{Global forest datasets}
\label{sec:org04e8818}

We considered two global or pantropical forest datasets to obtain maps of forest cover and forest cover change in New Caledonia: the Global Forest Change (GFC) dataset (\cslcitation{6}{Hansen et al. 2013}) and the Tropical Moist Forest (TMF) dataset (\cslcitation{10}{Vancutsem et al. 2021}). Both datasets rely on the analysis of the Landsat satellite image archive at a resolution of 30 m.

GFC is a global dataset providing estimates of (i) tree cover percentage for the year 2000, (ii) forest cover loss (change from forest to non-forest) on an annual basis from 2001 to 2025 (for GFC version 1.13 used in this study), and (iii) forest cover gain (canopy closure transitions from non-forest to forest) for the period 2000--2012 only. The tree cover percentage for the year 2000 provided by GFC ranges from 0 to 100, so that GFC can be used to characterize both dry and moist tropical forests. Nonetheless, a tree cover threshold must be chosen to define the initial forest cover for the year 2000. Some authors have used threshold values from 65 to 75 \% to define the moist forest in the tropics (\cslcitation{1}{Aleman et al. 2017}, \cslcitation{11}{Vieilledent et al. 2018}) but this threshold can be much lower for dry forest for example (\cslcitation{1}{Aleman et al. 2017}) and is variable between regions.

Contrary to GFC, the TMF dataset does not cover the entire globe but is available at the pantropical scale for evergreen forests, which can encompass both dry and moist tropical forests or only moist non-deciduous forests depending on the region considered. In New-Caledonia, a majority of the tree species from the three forest types (mangroves, dry and moist forests) are evergreen (\cslcitation{2}{BGCI 2021}), such that the TMF dataset can be used to characterize all forest types in New-Caledonia. TMF provides an annual product from the beginning of the archive (starting from 1990 depending on the location) to 2025 (for TMF version 2025 used in this study) where each pixel is categorized in 6 classes at each date: (1) undisturbed forest, (2) degraded forest, (3) deforested land, (4) forest regrowth, (5) permanent or seasonal water, and (6) other land cover.
\subsection{Wall-to-wall forest cover and forest cover change maps for New Caledonia}
\label{sec:org9aca491}
We first derived a wall-to-wall 2021 forest cover map for New-Caledonia using each of the two data-sets. We chose this year because we had access to very high resolution 2021 Pléiades satellite images on a large portion of New Caledonia land area for validating the forest cover map by photo-interpretation. To obtain the 2021 forest cover map for New Caledonia based on the GFC dataset, we first considered the forest pixels for the year 2000 assuming that a pixel was forested if it had a tree cover \(\geq \theta\). We tested values of \(\theta\) ranging from 10 to 90 in increments of 5. Second, we identified the deforested pixels for the years 2001 to 2020 using the forest cover loss product provided by GFC. We thus obtained a forest cover map for January 1\(^{\text{st}}\), 2021. We did not consider the forest gain to derive the 2021 forest cover map as it is only provided for the period 2000--2012 in the GFC dataset. We identified \(\theta^{\star}\), the optimal tree cover threshold for GFC which maximized the overall accuracy when comparing the 2021 forest cover map with the results from photo-interpretation (more details in section \ref{sec:org3145330}). For TMF, we considered the band ``Dec2020'' of the annual product which indicates the pixel state at the end of the year 2020 to obtain a forest cover map for January 1\(^{\text{st}}\), 2021. We classified pixels as forested if they fell in the following categories: undisturbed forest (1), degraded forest (2), or forest regrowth (4).

Second, we derived wall-to-wall forest cover change maps for the two periods 2004--2008 and 2008--2021. These years corresponded to years for which we had access to high resolution aerial images on a portion of New Caledonia for validating the forest cover change map by photo-interpretation. For GFC, we used the optimal tree cover threshold \(\theta^{\star}\) to define the forest in 2000 and derive the forest cover change maps for the two periods. For the first period (2004--2008), we defined the forest cover in 2004 as the set of forest pixels in 2000 which were not deforested in 2001--2003. Forest cover loss for that period was defined as pixels which were forest on January 1\(^{\text{st}}\), 2004 and deforested during the years 2004 to 2007. Forest cover in 2008 combined the forest pixels in 2000 which were not deforested in 2001--2007 and pixels of forest cover gain provided by GFC for the period 2000--2012 (assuming that a large part of this gain occurred during the period 2001--2007). For the second period (2008--2021), we did not considered forest cover gain which was not available after 2012 in the GFC dataset. Forest cover in 2008 was defined as forested pixels in 2000 which were not deforested in the period 2001--2007. Forest cover in 2021 was defined as forested pixels in 2000 which were not deforested in the period 2001--2020. Forest cover loss for the period 2008--2021 was obtained from the forest cover at the two dates. For TMF, we used the annual product to derive forest cover maps for the years 2004, 2008, and 2021. We classified pixels as forested if they fell in the following categories: undisturbed forest (1), degraded forest (2), or regrowth (4). For the two periods 2004--2008 and 2008--2021, and the two datasets GFC and TMF (except for GFC on 2008--2021), we ended up with four categories for forest cover change: ``loss'', ``gain'', ``stable forest'', and ``stable non-forest''. For GFC on 2008--2021, for which forest gain was unavailable, we only obtained three categories: ``loss'', ``stable forest'', and ``stable non-forest''.

Both GFC and TMF datasets are available on Google Earth Engine (GEE), a cloud-based platform by Google that combines a catalog of satellite imagery and environmental data with computing facilities (\cslcitation{5}{Gorelick et al. 2017}). Maps were obtained using the \texttt{geefcc} Python package (\url{https://pypi.org/project/geefcc/}) which allows computing and downloading forest cover change maps from the GFC and TMF datasets available on GEE.
\subsection{Photo-interpretation}
\label{sec:org2a76ad9}

\subsubsection{Photo-interpreting forest-cover for the year 2021}
\label{sec:org11dd216}

To obtain a reference dataset for the forest cover, we used very high spatial resolution (< 10 m) Pléiades satellite images covering the Southern Province of New Caledonia for the year 2021 (Fig. 1). Spatial resolution is 0.5 m for panchromatic images (one black and white band) and 2 m for multispectral images (four bands with red, green, blue and near infra-red bands). These non-free images were provided by the technical services of the Southern Province. The Southern Province covers about 38 \% of New Caledonia land area of 18,576 km\textsuperscript{2}.

We randomly sampled 1050 points in the Southern Province that were used to define the center of a cluster of nine (three by three) 30 m pixels. We obtained a total of 9450 sampled pixels. Three interpreters characterised the forest cover status of the pixels by photo-interpreting the Pléiades satellite images. They used the Collect Earth Online software (\url{https://www.collect.earth/}) to record their observations. Each interpreter classified a part (332) of the 1050 nine-pixel clusters independently from the other interpreters. Five percent of the clusters (54 clusters, corresponding to 486 pixels) were interpreted by the three interpreters in order to quantify the disagreement between interpreters (\cslcitation{9}{Olofsson et al. 2014}).

The interpreters assigned each pixel to one of three forest cover classes: ``forest'', ``non-forest'', or ``not interpretable''. The class ``not interpretable'' was applied to pixels obscured by clouds or shadows, or where the quality of the Pléiades imagery was insufficient to make a reliable determination. A pixel was classified as ``forest'' if a majority of the 30 by 30 m pixel was covered by forest. Forest was defined as a contiguous area of land covered by tree crowns of height greater than or equal to 5 m, exhibiting specific tonal, textural, and spatial patterns clearly distinct from other land cover types. The main other land cover types encountered in New Caledonia include bare soil, grasslands and savanna, mining maquis (a shrubland vegetation type characteristic of nutrient-poor ultramafic soils) and urban areas.

The interpreters made no distinction between different forest types existing in New Caledonia. Although moist tropical forest is the dominant forest type in New Caledonia (>500,000 ha, \cslcitation{8}{Myers et al. 2000}), relict dry tropical forests (10,000 ha) occur in the drier rain shadow of mountains along the west coast (\cslcitation{3}{Bouchet et al. 1995}), while tree plantations are found in certain specific locations (e.g. 4500 ha of \emph{Pinus caribaea} Morelet (Pinaceae) at the Tango plateau) and mangroves fringe the coastline (estimated at 29,627 ha in 2025 after \cslcitation{4}{Bunting et al. 2022}).

A particular challenge in New Caledonia lies in distinguishing mining maquis from forest on satellite imagery, as well-developed maquis stands can reach several meters in height, making them visually similar to low-stature forest formations, yet typically remaining below the 5 m height threshold used to define forest. Interpreters underwent a training phase prior to performing the pixel classification, during which difficult cases were discussed and clarification guidelines were established in order to minimize variability between interpreters.
\subsubsection{Photo-interpreting forest cover change for periods 2004--2008 and 2008--2021}
\label{sec:orgb9f651f}
To obtain reference datasets for the study of the forest cover change for the periods 2004--2008 and 2008--2021, we used a combination of 0.5 m resolution aerial photographs for the years 2003/2004, 2007/2008, and Pléiades satellite images for the year 2021. Aerial photographs were provided by the DITTT (``\emph{Direction des Infrastructures, de la Topographie et des Transports Terrestres}'') of the Government of New Caledonia. Intersection of aerial photographs in 2003/2004 and 2007/2008 covered an area of 109,355 ha, corresponding to 6 \% of New Caledonia land area. Intersection of aerial photographs in 2007/2008 and Pléiades images in 2021 covered an area of 558,077 ha, corresponding to 30 \% of New Caledonia land area.

Stratified sampling was employed to collect the reference datasets for forest cover change. This approach was adopted given the rarity of forest change events, which usually occur at a rate close to 1 \% per year (\cslcitation{9}{Olofsson et al. 2014}). Stratification was based on the classes of change obtained from the GFC and TMF datasets regarding the forest cover: ``gain'' or ``loss''. To assess change between 2004 and 2008, we drew a stratified sample of 1000 points: 600 in areas classified as loss by either GFC (300 points) or TMF (300 points), and 400 in areas classified as gain by either GFC (200 points) or TMF (200 points). Since GFC gain data are only available for the period 2000--2012 (\cslcitation{6}{Hansen et al. 2013}), we used this layer to sample the 200 points corresponding to GFC gain for the period 2004--2008. To assess change between 2008 and 2021, we drew a stratified sample of 850 points: 600 in areas classified as loss by either GFC (300 points) or TMF (300 points), and 250 in areas classified as gain by TMF. As before, points were used to define the center of a cluster of nine (three by three) 30 m pixels. We thus obtained a total of 9000 sampled pixels for the period 2004--2008 and 7650 pixels for the period 2008--2021.

Two interpreters shared the classification of the 9000 and 7650 pixels for the years 2004 and 2021 respectively, while three interpreters shared the classification of the 16650 pixels for the pivotal year 2008. As previously, five percent of the pixels were classified by several interpreters at each date to be able to estimate interpreter disagreement: by two interpreters for years 2004 and 2021, and by three interpreters for the year 2008. Pixels were classified as ``forest'', ``non-forest'', or ``not interpretable'' at each date using the same forest definition as before. From the forest cover classification at each date, we derived five classes of forest cover change for the two periods: ``loss'', ``gain'', ``stable forest'', ``stable non-forest'', and ``not interpretable''.
\subsection{Disagreement between interpreters}
\label{sec:orgd9e3561}

We obtained four datasets after photo-interpretation. One dataset, based on randomly sampled points, was used to analyze forest cover in 2021. The other three datasets, generated through stratified sampling, represented forest cover in 2004, 2008, and 2021 and were used to estimate forest cover change between these dates. Five percent of the sample clusters in each dataset were interpreted independently by two (2004, 2021) or three (2008, 2021 forest/non-forest) interpreters. To assess the level of disagreement among interpreters, each sample point interpreted by at least two of the three interpreters was analyzed across all four datasets. Raw disagreement was quantified as the proportion of sample points for which at least two interpreters assigned different land cover classes (``forest'', ``non-forest'', or ``not interpretable'').

We also quantified agreement with Krippendorff's alpha (\(\alpha_{K}\)) (\cslcitation{7}{Krippendorff 2008}), a chance-corrected measure of agreement, treating the data as nominal with three classes. Krippendorff's alpha was computed separately for each dataset and all datasets combined, using sample points interpreted by at least two interpreters. It is defined as \(\alpha_{K} = 1 - D_{o} / D_{e}\), where \(D_{o}\) is the observed disagreement, i.e. the proportion of discordant pairs of annotations within sample points (each point being weighted by its number of annotations), and \(D_{e}\) is the disagreement expected by chance, derived from the pooled class frequencies: \(D_{e} = (n^2 - \sum_k n_k^2) / (n(n-1))\), where \(n\) is the total number of annotations and \(n_k\) the number of annotations in class \(k\). Fore convenience, \(\alpha_{K}\), \(D_{e}\), and \(D_{o}\) were all expressed in percent in this study. A value of 100 for \(\alpha_{K}\) indicates perfect agreement and 0 indicates agreement no better than chance.

To assess whether one interpreter was less consistent than the others, we pooled the two datasets with three interpreters (2008 and 2021 forest/non-forest; 147 clusters, 1323 pixels) and applied two complementary approaches. First, we compared the alpha obtained with all three interpreters to the alpha obtained after removing each interpreter in turn, on the same set of pixels. A positive change in alpha after removing an interpreter indicates that this interpreter reduced the overall agreement. Confidence intervals (95\%) for the change in alpha were obtained by bootstrapping clusters (1,000 replicates), because the nine pixels of a cluster are not independent. Second, for pixels where two interpreters agreed and the third disagreed, we counted how often each interpreter was in the minority. Under equal reliability, each interpreter is expected to be in the minority in one third of cases, which we tested with a chi-squared goodness-of-fit test.
\subsection{Consensus reference datasets for forest cover and forest cover change}
\label{sec:orgbbcaf0c}

From the photo-interpretation results, we derived consensus reference datasets for forest cover in 2021 and for forest cover change over 2004--2008--2021. These datasets are the best available reference rather than an error-free ground truth (\cslcitation{9}{Olofsson et al. 2014}). Pixels classified as ``not interpretable'' could not be assessed and were excluded from the reference datasets. They represent a small proportion of the data (\(\leq\) 5.5\% across datasets).

For the 2021 forest cover dataset, a random subset of 486 pixels (5.1\% of the 9,450 pixels) was interpreted by three interpreters and assigned to one of three classes: ``forest'', ``non-forest'' or ``not interpretable''. The reference class of these pixels was the class chosen by the majority of interpreters. A majority was always reached, as no pixel received three different classes. Pixels for which the majority class was ``not interpretable'' (\(n = 60\)) were excluded. The remaining 95\% of pixels, interpreted by a single interpreter, kept the class assigned by that interpreter, except those classified as ``not interpretable'', which were also excluded (\(n = 461\)).

For the 2008 forest cover dataset used in the change analysis, the same majority rule was applied to the 837 pixels (5.0\% of 16,650) interpreted by three interpreters. A majority was always reached, and no pixel had ``not interpretable'' as majority class. Pixels interpreted by a single interpreter and classified as ``not interpretable'' were excluded (\(n = 65\)). For the 2004 and 2021 forest cover datasets used in the change analysis (including 9,000 and 7,650 pixels respectively), only two interpreters classified the replicated pixels (450 pixels in 2004 and 396 in 2021). Replicated pixels for which the two interpreters assigned different classes (96 and 102 pixels, respectively) were excluded, as no third opinion was available. Pixels classified as ``not interpretable'' by a single interpreter were also excluded (\(n = 154\) in 2004 and 156 in 2021).

For each of the two change periods (2004--2008 and 2008--2021), the change class of a pixel (``loss'', ``gain'', ``stable forest'' or ``stable non-forest'') was derived from its reference classes at the two dates. Pixels excluded at either date, because they were not interpretable or because no consensus was reached, were excluded from the change reference dataset (\(n = 263\) of 9000 for 2004--2008 and \(n = 294\) of 7650 for 2008--2021).
\subsection{Identifying the best GFC tree cover threshold for defining forest in New Caledonia}
\label{sec:org3145330}
We used the 2021 forest cover photo-interpretation dataset to estimate the optimal tree cover threshold \(\theta^{\star}\) for the Global Forest Change dataset to define forest in New Caledonia. We used this photo-interpretation dataset for two reasons: (i) its sample was uniform in space and adapted to a forest/non-forest classification, (ii) it covered the largest portion of land in New Caledonia (38 \%) compared with the other datasets used for analyzing forest cover-change. To determine \(\theta^{\star}\), we tested values of \(\theta\) ranging from 10 to 90 in increments of 5. For each value of \(\theta\), we derived a binary forest/non-forest classification for each pixel from the GFC tree cover in 2000 and tree cover loss layers from 2001 to 2020, and compared the resulting map against the photo-interpretation reference dataset by computing a confusion matrix. We computed the overall accuracy: \(\text{OA}(\theta) = (n_{11} + n_{00}) / (n_{11} + n_{10} + n_{01} + n_{00})\), where \(n_{11}\) and \(n_{00}\) are the number of true positives and true negatives, and \(n_{10}\) and \(n_{01}\) are the number of false positives and false negatives, respectively. The optimal threshold \(\theta^{\star}\) was defined as the value of \(\theta\) that maximized \(\text{OA}(\theta)\), reflecting the best agreement between the GFC-derived forest map and the photo-interpretation reference dataset. We used this optimal threshold \(\theta^{\star}\) to derive a forest cover map for year 2021 and forest cover change maps for 2004--2008 and 2008--2021 from the GFC dataset (see section \ref{sec:org9aca491}).
\subsection{Accuracy of forest cover and forest cover change maps}
\label{sec:orgb463663}

\subsubsection{Forest cover}
\label{sec:org142577d}

(OA, UA and PA. Should be weighted by the area).

We compared the TMF and GFC 2021 forest cover maps against the photo-interpretation reference dataset by computing confusion matrices. We computed the overall accuracy, as well as the user's and producer's accuracies of each class ``forest'' or ``non-forest''.  Olofsson et al. (\cslcitation{9}{2014}).
\subsubsection{Forest cover change}
\label{sec:org6782756}

(UA and commission errors; PA as an indication only. Not weighted by the area)

We computed the UA of the GFC and TMF maps for the classes ``loss'' and ``gain'' against the consensus reference. We also computed the PA of these classes as an indication only, since it is an upper bound (see the section on the effect of interpreter variability below). To characterize the commission errors, we computed the proportions of pixels of each consensus class (``loss'', ``gain'', ``stable forest'', ``stable non-forest'') among the pixels mapped as loss or gain by each product. These UA are sample-based indices, not weighted by the area of the strata (\cslcitation{9}{Olofsson et al. 2014}).
\subsection{Effect of interpreter variability on the accuracy of forest cover change maps}
\label{sec:org502859e}

The consensus reference datasets do not account for the uncertainty in the photo-interpretation. To quantify the effect of interpreter variability on the accuracy of forest cover change maps, we estimated the UA that would be obtained by a map identical to one interpreter's classification when evaluated against the classification of another interpreter. We approximated the first classification by the consensus, and described the second one by alternative reference datasets in which the reference class of each pixel was uncertain. This analysis was restricted to the two change periods (2004--2008 and 2008--2021), to the classes ``loss'' and ``gain'', and to the pixels retained in the consensus change reference datasets (\(n=33,300\)). We focused on the user's accuracy (UA), which is the proportion of pixels mapped in a change class that belong to the same class in the reference. We also computed the producer's accuracy (PA), which is the proportion of reference pixels of a change class that were mapped in the same class, as an indication only. Because pixels stable in both GFC and TMF were only sampled as neighbours of cluster centres, changes missed by both products were not representatively sampled, so that the PA is an upper bound (see section \ref{sec:orgb9f651f} for the sampling design of the photo-interpretation datasets).

For each pixel and each date, we defined the probability \(p\) that the reference class was ``forest''. For pixels interpreted by several interpreters (\(n=1625\)), \(p\) was the proportion of ``forest'' among the ``forest'' and ``non-forest'' labels (0, 0.5 or 1 if two interpreters, and 0, 1/3, 2/3 or 1 for three interpreters). For pixels interpreted by a single interpreter (\(n=31,282\)), \(p\) was the probability that another interpreter assigned the class ``forest'' given the label of the first interpreter. This probability was estimated separately for each date (2004, 2008, 2021) from all pairs of different interpreters on the replicated pixels, excluding pairs including a ``not interpretable'' label. This probability directly reflects the variability between interpreters, as it is estimated from the labels given by different interpreters to the same pixels. The alternative reference thus represents the reference that another interpreter could have produced, rather than the unknown true class. Assuming that interpretation errors at the two dates were independent, the probability that a pixel was a loss in the alternative reference was \(p_1 (1 - p_2)\), and the probability that it was a gain was \((1 - p_1) p_2\), where \(p_1\) and \(p_2\) are the probabilities of forest at the first and second date of the period.

We then estimated the maximum UA and PA that a map can reach given interpreter variability (hereafter the ceilings). For this, we considered a hypothetical map identical to the consensus change classes. Because this map is fixed, the set of pixels mapped in each class is known, and the expected UA of the classes ``loss'' and ``gain'' is the mean of the above probabilities over the pixels of each class. This is the value that would be obtained by averaging the UA over an infinite number of simulated alternative reference datasets, so no random draws were needed. The ceiling of the PA was computed in the same way, as the expected number of pixels of the class in the alternative reference that are mapped in this class by the hypothetical map, divided by the expected total number of pixels of the class in the alternative reference (ratio of expectations, which approximates the expectation of the ratio given the sample sizes). We then compared the UA and PA obtained for the GFC and TMF maps for the classes ``loss'' and ``gain'' with their ceilings.
\subsection{Comparison with the 2015 forest cover map derived from aerial photographs}
\label{sec:orgd64eb3c}

\subsection{Deriving 2005--2015--2025 forest cover change maps from TMF for New Caledonia}
\label{sec:org5e3cd76}

\subsection{Analysis of forest fragmentation for the year 2025}
\label{sec:org1ff020b}

\newpage
\section{Results}
\label{sec:org360e2a1}

\subsection{Disagreement between interpreters}
\label{sec:orgbd807f9}

Interpreter disagreement was substantial (Table \ref{tab:orga9cc952}). Among the pixels interpreted by at least two interpreters (about 5\% of the sample, 2,169 pixels in total), 21.3 to 25.8\% were assigned different classes by at least two interpreters, i.e. 24.0\% overall. Chance-corrected agreement was only moderate to substantial, with Krippendorff's alpha ranging from 52.1\% (2004) to 72.7\% (2021 forest/non-forest). There was no evidence that a single interpreter drove this disagreement. Removing any one interpreter changed alpha by only -2.6 to +1.4 percentage points, and all bootstrap 95\% CIs included zero (Table \ref{tab:orgcfc36d6}). Among the 322 pixels where two interpreters agreed and the third disagreed, the three interpreters were in the minority in similar proportions (A: 28.6\%, O: 35.7\%, S: 35.7\%). These proportions did not differ significantly from the 33.3\% expected under equal reliability (\(\chi^2\) = 3.29, df = 2, p = 0.19). Disagreement was globally lower for the 2021 forest/non-forest dataset (D\_o = 15.6\%) than for the three other datasets (16.6\% for 2008, 21.3\% for 2004 and 25.8\% for 2021), consistent with its higher alpha.
\subsection{Accuracy of forest cover maps}
\label{sec:org2270d9b}

\subsection{Accuracy of forest cover change maps}
\label{sec:org1556643}

The consensus change reference datasets included 8,737 pixels for 2004--2008 and 7,327 pixels for 2008--2021. Of these, 1,037 and 1,120 pixels were classified as loss, and 1,133 and 698 pixels as gain, for 2004--2008 and 2008--2021 respectively. Given interpreter variability, the ceiling for the UA of the loss class, i.e. the expected UA of a map identical to the consensus, was 0.60 for 2004--2008 and 0.62 for 2008--2021. The ceiling for the UA of the gain class was 0.68 for 2004--2008 and 0.64 for 2008--2021 (Table \ref{tab:orge0d9ad2}).

Against the consensus reference, the UA of the loss class was 0.48 for GFC (630 pixels mapped as loss) and 0.35 for TMF (445 pixels) for 2004--2008, and 0.33 for GFC (531 pixels) and 0.34 for TMF (1,113 pixels) for 2008--2021. These values correspond to 54 to 80\% of the ceiling. The UA of the gain class was 0.22 for GFC (373 pixels) and 0.24 for TMF (227 pixels) for 2004--2008, i.e. 32 and 35\% of the ceiling, and 0.32 for TMF (464 pixels) for 2008--2021, i.e. 50\% of the ceiling. GFC gain was not evaluated for 2008--2021, as it is not available after 2012.

Among the pixels mapped as loss, the main source of commission error was pixels classified as stable non-forest in the consensus reference (39 to 47\%), far ahead of stable forest (7 to 17\%) and gain (4 to 8\%) (Table \ref{tab:orge0d9ad2}). The composition of the pixels mapped as gain was more variable between products and periods, mostly between stable non-forest and stable forest. For 2004--2008, 66\% (GFC) and 31\% (TMF) of the pixels mapped as gain were stable non-forest, while 9\% (GFC) and 41\% (TMF) were stable forest. For 2008--2021, 47\% of the pixels mapped as gain by TMF were stable non-forest and 15\% were stable forest. Only 3 to 7\% of the pixels mapped as gain were loss in the reference.

As an indication, the PA of the loss class was 0.29 for GFC and 0.15 for TMF for 2004--2008, and 0.16 for GFC and 0.34 for TMF for 2008--2021, for ceilings of 0.45 and 0.49 respectively (i.e. 32 to 70\% of the ceiling). The PA of the gain class was 0.07 for GFC and 0.05 for TMF for 2004--2008 (ceiling 0.44, i.e. 16 and 11\%), and 0.21 for TMF for 2008--2021 (ceiling 0.31, i.e. 69\%). For all products, periods and classes, the PA was lower than or equal to the UA (Table \ref{tab:orge0d9ad2}).
\subsection{Forest cover, deforestation, regrowth and forest fragmentation in 2005--2025}
\label{sec:org5d37d2a}

\newpage
\section{Discussion}
\label{sec:org99a38b8}

\subsection{Effect of interpreter variability on the accuracy of forest cover change maps}
\label{sec:org57e41a7}

User's accuracies for forest loss were low (0.33--0.48 depending on the product and period; Table \ref{tab:orge0d9ad2}). Such values must be interpreted in the light of the uncertainty of the reference data. Interpreters disagreed on 21--26 \% of the pixels (Table \ref{tab:orga9cc952}), and forest loss requires two correct interpretations (forest at the first date, non-forest at the second). A map identical to the consensus would thus reach a user's accuracy of only 0.60--0.62 for loss when evaluated against simulated reference data-sets reflecting interpreter variability. Given the difficulty of the exercise, the two products reached more than half of this ceiling. Most of the remaining commission errors were pixels classified as stable non-forest by the interpreters (39--47 \% of the pixels mapped as loss), which suggests a confusion between forest and other tree-covered or shrubby vegetation types (notably mining maquis) in the products. Forest gain was even harder to detect (user's accuracy of 0.22--0.32), with 9--41 \% of the pixels mapped as gain being stable forest and 31--66 \% stable non-forest according to the interpreters. Producer's accuracies, given as an indication only, were low (0.15--0.34 for loss and 0.05--0.21 for gain) and must be considered as upper bounds. The sample was stratified on the loss and gain classes of GFC and TMF, and only the central pixel of each cluster was drawn from these strata. Pixels that are stable in both products were thus only sampled as neighbouring pixels of cluster centres, so changes missed by both products were not representatively sampled. Low accuracies for change should therefore be expected, and the figures reported here are sample-based indices rather than area-weighted estimates (\cslcitation{9}{Olofsson et al. 2014}). The high commission errors observed here do not show that forest loss is overestimated in net terms, as omission errors may partly compensate for them. The indicative producer's accuracies, which are upper bounds, were never higher than the user's accuracies, so the data give no sign of a net overestimation of forest loss. Only an area-weighted estimation could confirm this. A natural perspective would be to complement the current stratified samples with a random sample of the stratum of pixels stable in both products, within the area covered by both the 2007/2008 aerial photographs and the 2021 Pléiades images. Such a design, with mutually exclusive strata and known inclusion probabilities, would allow estimating unbiased producer's accuracies, omission errors, and area-weighted accuracies and change areas with their confidence intervals (\cslcitation{9}{Olofsson et al. 2014}).
\subsection{Figures of forest cover and deforestation in New-Caledonia}
\label{sec:org1f5159b}

\subsection{Fragmentation of New-Caledonian forests}
\label{sec:org54add5a}

\subsection{Impact for biodiversity conservation}
\label{sec:org6aba26d}

\newpage
\section{References}
\label{sec:org4ee791d}

\begin{cslbibliography}{1}{0}
\cslbibitem{1}{Aleman, J. C., O. Blarquez, S. Gourlet-Fleury, L. Bremond, and C. Favier. 2017. \href{https://doi.org/10.1038/srep41393}{Tree cover in central africa: determinants and sensitivity under contrasted scenarios of global change}. Scientific Reports 7.}

\cslbibitem{2}{BGCI. 2021. State of the world’s trees. Botanic Gardens Conservation International.}

\cslbibitem{3}{Bouchet, P., T. Jaffre, and J.-M. Veillon. 1995. \href{https://doi.org/10.1007/bf00058425}{Plant extinction in new caledonia: protection of sclerophyll forests urgently needed}. Biodiversity and Conservation 4:415–428.}

\cslbibitem{4}{Bunting, P., A. Rosenqvist, L. Hilarides, R. M. Lucas, N. Thomas, T. Tadono, T. A. Worthington, M. Spalding, N. J. Murray, and L.-M. Rebelo. 2022. \href{https://doi.org/10.3390/rs14153657}{Global mangrove extent change 1996–2020: Global mangrove watch version 3.0}. Remote sensing 14:3657.}

\cslbibitem{5}{Gorelick, N., M. Hancher, M. Dixon, S. Ilyushchenko, D. Thau, and R. Moore. 2017. \href{https://doi.org/10.1016/j.rse.2017.06.031}{Google earth engine: Planetary-scale geospatial analysis for everyone}. Remote Sensing of Environment 202:18–27.}

\cslbibitem{6}{Hansen, M. C., P. V. Potapov, R. Moore, M. Hancher, S. A. Turubanova, A. Tyukavina, D. Thau, S. V. Stehman, S. J. Goetz, T. R. Loveland, A. Kommareddy, A. Egorov, L. Chini, C. O. Justice, and J. R. G. Townshend. 2013. \href{https://doi.org/10.1126/science.1244693}{High-resolution global maps of 21st-century forest cover change}. Science 342:850–853.}

\cslbibitem{7}{Krippendorff, K. 2008. \href{https://doi.org/10.1080/19312450802467134}{Systematic and random disagreement and the reliability of nominal data}. Communication methods and measures 2:323–338.}

\cslbibitem{8}{Myers, N., R. A. Mittermeier, C. G. Mittermeier, G. A. B. da Fonseca, and J. Kent. 2000. Biodiversity hotspots for conservation priorities. Nature 403:853–858.}

\cslbibitem{9}{Olofsson, P., G. M. Foody, M. Herold, S. V. Stehman, C. E. Woodcock, and M. A. Wulder. 2014. \href{https://doi.org/10.1016/j.rse.2014.02.015}{Good practices for estimating area and assessing accuracy of land change}. Remote Sensing of Environment 148:42–57.}

\cslbibitem{10}{Vancutsem, C., F. Achard, J.-F. Pekel, G. Vieilledent, S. Carboni, D. Simonetti, J. Gallego, L. E. O. C. Aragão, and R. Nasi. 2021. \href{https://doi.org/10.1126/sciadv.abe1603}{Long-term (1990–2019) monitoring of forest cover changes in the humid tropics}. Science Advances 7.}

\cslbibitem{11}{Vieilledent, G., C. Grinand, F. A. Rakotomalala, R. Ranaivosoa, J.-R. Rakotoarijaona, T. F. Allnutt, and F. Achard. 2018. \href{https://doi.org/https://doi.org/10.1016/j.biocon.2018.04.008}{Combining global tree cover loss data with historical national forest cover maps to look at six decades of deforestation and forest fragmentation in madagascar}. Biological Conservation 222:189–197.}

\end{cslbibliography}
\section{Figures}
\label{sec:orgb1d02e2}

\subsection{Figure X}
\label{sec:orgf529d91}

\begin{figure}[H]
\centering
\includegraphics[width=\linewidth]{figures/fcc_tmf_10yr.png}
\caption{\label{fig:org4c467ad}\textbf{Map of forest cover change for 2015--2025.}}
\end{figure}

\newpage
\subsection{Figure X}
\label{sec:org208919e}

\begin{figure}[H]
\centering
\includegraphics[width=0.8\linewidth]{figures/distance_edge_histogram.png}
\caption{\label{fig:org64c87ef}\textbf{Distribution of 30 m forest pixels in classes of distance to forest edge.} For year 2025.}
\end{figure}

\newpage
\subsection{Figure X}
\label{sec:orgb37e82b}

\begin{figure}[H]
\centering
\includegraphics[width=\linewidth]{figures/comp_amap_tmf.png}
\caption{\label{fig:org30158fd}\textbf{Comparing TMF derived map with forest cover map derived from aerial photographs for the year 2015.}}
\end{figure}

\newpage
\subsection{Figure X}
\label{sec:org6abd25b}

\begin{figure}[H]
\centering
\includegraphics[width=0.8\linewidth]{figures/plot_accuracy_gfc_perc.pdf}
\caption{\label{fig:org96342c7}\textbf{Identifying the best tree cover threshold for defining forest with the GFC product.}}
\end{figure}

\newpage
\section{Tables}
\label{sec:orgc0d9609}

\begin{table}[htbp]
\caption{\label{tab:orga9cc952}\textbf{Disagreement statistics per dataset.} This table presents the raw percentage of disagreement \(\%_{\text{disag}}\) and the Krippendorff's alpha (\(\alpha_{K}\)) for each dataset and for all datasets combined. The dataset \texttt{for2021\_fnf} was used to evaluate the forest/non-forest cover in 2021. The three other datasets \texttt{for2000}, \texttt{for2008}, and \texttt{for2021} were used to evaluate the forest cover change in 2000--2008--2021. Krippendorff's alpha (here in \%) ranges from 100 (perfect agreement) to 0 (agreement no better than chance). \(n_{C}\) is the number of pixel clusters (3 \texttimes{} 3 pixels with 30 m resolution each), \(n_{p}\) is the number of pixels, \(n_{\text{rep}}\) is the number of pixels analyzed by several interpreters, \(\%_{\text{rep}}\) is the percentage of pixels analyzed by several interpreters, \(n_{\text{disag}}\) is the number of disagreements, \(D_{o}\) is the observed disagreement (percentage of discordant pairs among all pairs), \(D_{e}\) is the expected percentage of disagreement by chance. \(D_o\) differs from the raw percentage of disagreement, which only counts the proportion of pixels with disagreement.\label{tab:disag}\\
}
\centering
\small
\begin{tabular}{lrrrrrrrrr}
\toprule
Dataset & \(n_{C}\) & \(n_{p}\) & \(n_{\text{rep}}\) & \(\%_{\text{rep}}\) & \(n_{\text{disag}}\) & \(\%_{\text{disag}}\) & \(D_{o}\) & \(D_{e}\) & \(\alpha_{K}\)\\
\midrule
for2004 & 1000 & 9000 & 450 & 5.0 & 96 & 21.3 & 21.3 & 44.5 & 52.08\\
for2008 & 1850 & 16650 & 837 & 5.0 & 208 & 24.9 & 16.6 & 46.5 & 64.40\\
for2021 & 850 & 7650 & 396 & 5.2 & 102 & 25.8 & 25.8 & 55.2 & 53.36\\
for2021\_fnf & 1050 & 9450 & 486 & 5.1 & 114 & 23.5 & 15.6 & 57.2 & 72.65\\
All combined & 4750 & 42750 & 2169 & 5.1 & 520 & 24.0 & 18.4 & 53.3 & 65.53\\
\bottomrule
\end{tabular}
\end{table}

\begin{table}[htbp]
\caption{\label{tab:orge0d9ad2}\textbf{User's and producer's accuracies of the forest cover change maps and composition of the pixels mapped as loss or gain.} For each period, product (GFC or TMF) and mapped class (loss or gain), the table gives the number of pixels mapped in the class (\(n_{\text{map}}\)), the user's accuracy (UA) and the producer's accuracy (PA) against the consensus reference, the ceilings (Ceil.) of the UA and PA, and the percentage of pixels of each consensus reference class (loss, gain, stable forest SF, stable non-forest SNF) among the pixels mapped in the class. The ceiling is the expected UA of a map identical to the consensus when evaluated against alternative reference datasets accounting for interpreter variability. The ceiling of the PA is the expected PA of the same map. GFC gain was not available for 2008--2021. Values are sample-based indices, not weighted by the area of the strata. The PA is given as an indication only and is an upper bound, as changes missed by both products were not representatively sampled.\label{tab:change-ua}\\
}
\centering
\small
\begin{tabular}{rllrrrrrrrrr}
\toprule
Period & Product & Class & \(n_{\text{map}}\) & UA & UA & PA & PA & Loss & Gain & SF & SNF\\
 &  &  &  &  & Ceil. &  & Ceil. & (\%) & (\%) & (\%) & (\%)\\
\midrule
2004--2008 & GFC & loss & 630 & 0.48 & 0.60 & 0.29 & 0.45 & 48.4 & 5.6 & 7.0 & 39.0\\
2004--2008 & GFC & gain & 373 & 0.22 & 0.68 & 0.07 & 0.44 & 3.2 & 21.7 & 9.1 & 66.0\\
2004--2008 & TMF & loss & 445 & 0.35 & 0.60 & 0.15 & 0.45 & 35.3 & 7.9 & 16.6 & 40.2\\
2004--2008 & TMF & gain & 227 & 0.24 & 0.68 & 0.05 & 0.44 & 4.4 & 24.2 & 40.5 & 30.8\\
2008--2021 & GFC & loss & 531 & 0.33 & 0.62 & 0.16 & 0.49 & 33.3 & 5.1 & 14.5 & 47.1\\
2008--2021 & TMF & loss & 1113 & 0.34 & 0.62 & 0.34 & 0.49 & 34.2 & 4.0 & 16.4 & 45.4\\
2008--2021 & TMF & gain & 464 & 0.32 & 0.64 & 0.21 & 0.31 & 6.5 & 31.9 & 15.1 & 46.6\\
\bottomrule
\end{tabular}
\end{table}

\newpage
\section{Acknowledgements}
\label{sec:org7121ed2}

This research received fundings from the SURFOR project funded by DAFE (``\emph{Direction du service d’État de l’Agriculture, de la Forêt et de l'Environnement}''). The authors declare no conflicts of interest.
\section{Authors contribution}
\label{sec:orga07eac6}

GV, AB and FA conceived and designed the study. FA and AB secured the research funding. AB and OB, assisted by a graduate trainee, carried out the photo-interpretation. GV, with support from LL and MD, performed the data analysis. GV drafted the manuscript, with significant input from LL. PB, TI, and NE contributed to the derivation of the 2015 forest-cover map from aerial photographs. PH and DB provided critical review of the final forest-cover change map. All authors read, reviewed, and approved the final version of the manuscript.
\section{Data and code availability}
\label{sec:org65c2c74}

The study's raw data, code and outputs are available on \href{https://github.com/ghislainv/fcc-nc-surfor-v2}{GitHub} and will be permanently archived in the Cirad Dataverse as soon as the article is accepted for publication.
\section{Supplementary information}
\label{sec:orgf1ed016}

\renewcommand{\thetable}{S\arabic{table}}
\renewcommand{\theHtable}{S\arabic{table}}
\setcounter{table}{0}

\begin{table}[htbp]
\caption{\label{tab:orgcfc36d6}\textbf{Effect of removing each interpreter on Krippendorff's alpha.} Analysis based on the two datasets interpreted by three interpreters (\texttt{for2008} and \texttt{for2021\_fnf}; 147 clusters, 1323 pixels). Interpreters are identified by their initials (A, O, S). \(\Delta\alpha_{K}\) is the change in \(\alpha_{K}\) (in percentage points) after removing the interpreter, relative to the \(\alpha_{K}\) obtained with the three interpreters on the same pixels. A positive value indicates that the interpreter reduced the overall agreement. The 95\% confidence interval (CI) of \(\Delta\alpha_{K}\) was obtained by bootstrapping clusters of pixels (1,000 replicates).\label{tab:loo}\\
}
\centering
\small
\begin{tabular}{lrrr}
\toprule
Interpreter & \(\Delta\alpha_{K}\) & CI 2.5\% & CI 97.5\%\\
\midrule
S & 1.4 & -3.0 & 6.3\\
A & -2.6 & -7.5 & 2.3\\
O & 0.9 & -3.5 & 5.5\\
\bottomrule
\end{tabular}
\end{table}
\end{document}
